\section{Περιγραφή της εργασίας}



Για να εντοπίσουμε την θέση μίας πηγής ρύπου στο αέρα εσωτερικού χώρου,
αναπτύξαμε ένα δίκτυο αισθητήρων μέσω IoT για την παρακολούθηση του αέρα και την αποθήκευση των τιμών των αισθητήρων,όπως και τον χρηστών.
Πραγματοποίησαμε πολλαπλά πειράματα όπου εισάγαμε την πηγή ρύπου στο χώρο, με σκοπό την χρήστη τους στην ανάπτυξη των μηχανικών μοντέλων. 
Σαν μηχανικά μοντέλα χρησιμοποιήσαμε μοντέλα παλινδρόμισης για να βρούμε την απόσταση της πηγής από τον κάθε αισθητήρα,
με την μετέπειτα χρήση τριαπλευρικού εντοπισμού για την εύρεση της τελικής θέσης.

Θα ξεκινήσουμε με το κομμάτι που αφορά τους αισθητήρες,
περιγράφοντας πρώτα τα περιβάλλοντα προγραμματισμού, framework και βιβλιοθήκες που χρησιμοποιήθηκαν για τον προγραμματισμό των  μικροελεκτών που διαχειρίζεται τον αισθητήρα και αποστέλνει τα δεδομένα στον διακομιστή (server).
Περιγράφουμε τον μικροελετκή ESP32 και τον αισθητήρα BME680,όπως και το πως προγραμματιστήκε και τους λόγους αυτού του προγραμματισμού.

Συνεχίζουμε αναλύοντας τις τεχνολογίες που χρησιμοποιήσαμε για την ανάπτυξη του διαμοσιστή,της βάσης δεδομένων που αποθηκεύονται τα δεδομένα όπως και την οπτικοποίηση σε ζωντανό χρόνο.

Περιγράφουμε τις λεπτομέρειες των συνθηκών διατέλεσης των πειραμάτων και αποφάσης που πάρθηκαν κατά την διάρκεια, όπως και μια αναλυτική περιγραφή του τι γινόταν σε κάθε πείραμα.

Έπειτα,παραθέτουμε τις βιβλιοθήκες και το περιβάλλον που χρησιμοποιήσαμε για την υλοποίηση του κώδικα σχετικά με την ανάλυση την επεξεργασία και ανάλυση των δεδομένων των πειραμάτνω,όπως και των μηχανικών μοντέλων.
Περιγράφουμε αναλυτικά τον κώδικα που υλποποιήσαμε στο τελευταίο κομμάτι,που χωρίζεται σε 4 κομμάτσια:
Στην αίτηση και απόκτηση των δεδομένων από τον server,την επεξεργασία τους, την άντληση πληροφορίων μέσω οπτικοποίησης και τέλος την δημιουργία μηχανικών μοντέλων.